\documentclass[aspectratio=169]{beamer}

\usetheme{Madrid}
\usecolortheme{default}

\usepackage{amsmath, amssymb, bm}
\usepackage{booktabs}

\title[Week 4 Reading]{Week 4 Reading: From LQR to Trajectory Optimization}
\subtitle{Local Stabilization, Nominal Trajectory, MPC, and iLQR}
\author{Nguyen Dang Phuc}
\date{}

\begin{document}

% Slide 1
\begin{frame}
  \titlepage
\end{frame}

% Slide 2
\begin{frame}{Role of LQR in robotics}
\large

LQR is very useful in robotics, but mainly as a \textbf{local stabilizer}.

\vspace{0.4cm}

\begin{itemize}
    \item It stabilizes the system around an equilibrium point or a nominal trajectory.
    \item It gives an interpretable feedback controller:
    \[
    \bm{u} = \bm{u}_0 - K(\bm{x}-\bm{x}_0)
    \]
    \item It is not designed to solve global motion generation.
\end{itemize}

\vspace{0.4cm}

\begin{block}{Key distinction}
LQR answers: \textit{How do we stabilize around a known point or trajectory?}

It does not answer: \textit{How do we generate the trajectory in the first place?}
\end{block}

\end{frame}

% Slide 3
\begin{frame}{Optimal control problem}
\large

Consider a nonlinear dynamical system:

\[
\dot{\bm{x}} = f(\bm{x}, \bm{u})
\]

We want to choose a control trajectory that minimizes:

\[
J =
\ell_f(\bm{x}(t_f))
+
\int_{t_0}^{t_f}
\ell(\bm{x}(t), \bm{u}(t))\,dt
\]

subject to:

\[
\dot{\bm{x}}(t) = f(\bm{x}(t), \bm{u}(t))
\]

\vspace{0.3cm}

\begin{block}{Interpretation}
Optimal control chooses actions while respecting the physical dynamics.
\end{block}

\end{frame}

% Slide 4
\begin{frame}{The LQR problem}
\large

LQR is a special optimal control problem.

\vspace{0.2cm}

Linear dynamics:

\[
\dot{\bm{x}} = A\bm{x} + B\bm{u}
\]

Quadratic infinite-horizon cost:

\[
J =
\int_0^\infty
\left(
\bm{x}^T Q \bm{x}
+
\bm{u}^T R \bm{u}
\right)dt
\]

where:

\[
Q = Q^T \succeq 0,
\qquad
R = R^T \succ 0
\]

\vspace{0.2cm}

\begin{itemize}
    \item $Q$: penalty on state error.
    \item $R$: penalty on control effort.
\end{itemize}

\end{frame}
% Slide 5
\begin{frame}[t]{From HJB to the optimal LQR policy}
\large

Infinite-horizon HJB equation:

\[
0 =
\min_{\bm{u}}
\left[
\ell(\bm{x},\bm{u})
+
\nabla_{\bm{x}} V(\bm{x})^{T}
f(\bm{x},\bm{u})
\right]
\]

\vspace{-0.1cm}

For the LQR problem:

\[
\begin{aligned}
f(\bm{x},\bm{u}) 
&= A\bm{x}+B\bm{u}, \\[0.15cm]
\ell(\bm{x},\bm{u})
&=
\bm{x}^TQ\bm{x}
+
\bm{u}^TR\bm{u}.
\end{aligned}
\]

\vspace{-0.1cm}

Quadratic value function:
\[
V(\bm{x}) = \bm{x}^T S \bm{x},
\qquad
S=S^T \succeq 0
\]
\end{frame}
% Slide 6
\begin{frame}{Deriving the Riccati equation}
\large

Substitute into HJB:

\[
0 =
\min_{\bm{u}}
\left[
\bm{x}^TQ\bm{x}
+
\bm{u}^TR\bm{u}
+
2\bm{x}^T S(A\bm{x}+B\bm{u})
\right]
\]

Define:

\[
F(\bm{u})
=
\bm{x}^TQ\bm{x}
+
\bm{u}^TR\bm{u}
+
2\bm{x}^TSA\bm{x}
+
2\bm{x}^TSB\bm{u}
\]

Take derivative w.r.t. $\bm{u}$:

\[
\frac{\partial F}{\partial \bm{u}}
=
2R\bm{u}
+
2B^TS\bm{x}
=
0
\]

Therefore:

\[
\bm{u}^*
=
-R^{-1}B^TS\bm{x}
=
-K\bm{x}
\]

\end{frame}

% Slide 7
\begin{frame}{Solving Riccati gives the feedback gain}
\large

Substitute the optimal policy back into HJB:

\[
\bm{u}^* = -R^{-1}B^TS\bm{x}
\]

This yields:

\[
0 =
\bm{x}^T
\left[
A^TS + SA
-
SBR^{-1}B^TS
+
Q
\right]
\bm{x}
\]

Since this must hold for all $\bm{x}$:

\[
\boxed{
A^TS + SA
-
SBR^{-1}B^TS
+
Q
=
0
}
\]

This is the \textbf{algebraic Riccati equation}.

\vspace{0.3cm}

Once $S$ is solved:

\[
K = R^{-1}B^TS
\]

\[
\bm{u}^* = -K\bm{x}
\]

\end{frame}

% Slide 8
\begin{frame}{Why LQR is still useful for nonlinear systems}
\large

Most robotic systems are nonlinear:

\[
\dot{\bm{x}} = f(\bm{x}, \bm{u})
\]

Around a stabilizable operating point $(\bm{x}_0,\bm{u}_0)$:

\[
f(\bm{x}_0,\bm{u}_0)=0
\]

Define error coordinates:

\[
\bar{\bm{x}} = \bm{x}-\bm{x}_0,
\qquad
\bar{\bm{u}} = \bm{u}-\bm{u}_0
\]

Linearize:

\[
\dot{\bar{\bm{x}}}
\approx
A\bar{\bm{x}} + B\bar{\bm{u}}
\]

Then LQR gives:

\[
\bm{u}^*
=
\bm{u}_0
-
K(\bm{x}-\bm{x}_0)
\]

\end{frame}

% Slide 9
\begin{frame}{The limitation of local control}
\large

LQR is powerful, but local.

\vspace{0.3cm}

It works when:

\[
\bm{x} \approx \bm{x}_0
\]

or when the system stays close to a known trajectory.

\vspace{0.4cm}

But many robotic tasks require:

\[
\text{initial state}
\longrightarrow
\text{desired behavior}
\]

through large nonlinear motion.

\vspace{0.4cm}

\begin{block}{Limitation}
LQR can stabilize around a target, but it does not generate the global motion needed to reach that target.
\end{block}

\end{frame}

% Slide 10
\begin{frame}{If a nominal trajectory is already given}
\large

Assume we already have a feasible nominal trajectory:

\[
\bm{x}_0(t), \bm{u}_0(t)
\]

satisfying:

\[
\dot{\bm{x}}_0(t)
=
f(\bm{x}_0(t),\bm{u}_0(t))
\]

Define trajectory-relative errors:

\[
\bar{\bm{x}}(t)
=
\bm{x}(t)-\bm{x}_0(t)
\]

\[
\bar{\bm{u}}(t)
=
\bm{u}(t)-\bm{u}_0(t)
\]

Linearize along the trajectory:

\[
\dot{\bar{\bm{x}}}
\approx
A(t)\bar{\bm{x}}
+
B(t)\bar{\bm{u}}
\]

\end{frame}

% Slide 11
\begin{frame}{From nominal trajectory to feedback}
\large

Using time-varying LQR, we obtain:

\[
\bar{\bm{u}}^*(t)
=
-K(t)\bar{\bm{x}}(t)
\]

Therefore:

\[
\boxed{
\bm{u}^*(t)
=
\bm{u}_0(t)
-
K(t)
\left(
\bm{x}(t)-\bm{x}_0(t)
\right)
}
\]

\vspace{0.4cm}

\begin{block}{Interpretation}
The nominal trajectory gives the planned motion.  
Time-varying LQR gives local correction around that motion.
\end{block}

\vspace{0.3cm}

\[
\text{open-loop trajectory}
+
\text{local feedback}
=
\text{locally stabilized motion}
\]

\end{frame}

% Slide 12
\begin{frame}{Trajectory optimization}
\large

Now we return to the missing question:

\[
\textbf{How do we find the nominal trajectory?}
\]

\vspace{0.3cm}

Trajectory optimization searches for:

\[
\bm{x}_0(t), \bm{u}_0(t)
\]

by solving:

\[
\min_{\bm{u}(\cdot)}
\ell_f(\bm{x}(t_f))
+
\int_{t_0}^{t_f}
\ell(\bm{x}(t),\bm{u}(t))dt
\]

subject to:

\[
\dot{\bm{x}}(t)
=
f(\bm{x}(t),\bm{u}(t))
\]

\[
\bm{x}(t_0)=\bm{x}_0
\]

\end{frame}

% Slide 13
\begin{frame}{How to model trajectory optimization}
\large

The continuous constraint:

\[
\dot{\bm{x}}(t)
=
f(\bm{x}(t),\bm{u}(t)),
\qquad
\forall t
\]

is difficult to enforce directly.

\vspace{0.4cm}

So trajectory optimization methods replace it with discrete approximations.

\vspace{0.3cm}

Three common formulations:

\begin{enumerate}
    \item Direct transcription
    \item Direct shooting
    \item Direct collocation
\end{enumerate}

\end{frame}

% Slide 14
\begin{frame}{Direct transcription}

Decision variables:

\[
\bm{x}[0],\ldots,\bm{x}[N],
\qquad
\bm{u}[0],\ldots,\bm{u}[N-1]
\]

Optimization problem:

\[
\min_{\bm{x}[\cdot],\bm{u}[\cdot]}
\ell_f(\bm{x}[N])
+
\sum_{n=0}^{N-1}
\ell(\bm{x}[n],\bm{u}[n])
\]

subject to discrete dynamics:

\[
\bm{x}[n+1]
=
f_d(\bm{x}[n],\bm{u}[n])
\]

\[
\bm{x}[0]=\bm{x}_0
\]

\vspace{0.3cm}


\end{frame}

% Slide 15
\begin{frame}{Direct shooting}
\large

Decision variables:

\[
\bm{u}[0],\ldots,\bm{u}[N-1]
\]

States are generated by forward simulation:

\[
\bm{x}[n+1]
=
f_d(\bm{x}[n],\bm{u}[n])
\]

So:

\[
\bm{x}[n]
=
\phi_n(\bm{x}_0,\bm{u}[0],\ldots,\bm{u}[n-1])
\]

The optimization becomes:

\[
\min_{\bm{u}[\cdot]}
\ell_f(\bm{x}[N])
+
\sum_{n=0}^{N-1}
\ell(\bm{x}[n],\bm{u}[n])
\]

\vspace{0.3cm}

\begin{block}{Intuition}
Optimize only controls; states are implicit through rollout.
\end{block}

\end{frame}

% Slide 16
\begin{frame}{Direct collocation}
\large

Direct collocation represents states and controls as piecewise polynomials.

\vspace{0.3cm}

Decision variables are values at break points:

\[
\bm{x}[0],\ldots,\bm{x}[N],
\qquad
\bm{u}[0],\ldots,\bm{u}[N]
\]

Between two break points, choose a collocation point:

\[
t_{c,k}
=
\frac{1}{2}(t_k+t_{k+1})
\]

Then enforce dynamics at that point:

\[
\dot{\bm{x}}(t_{c,k})
=
f(\bm{x}(t_{c,k}),\bm{u}(t_{c,k}))
\]

\end{frame}

% Slide 17
\begin{frame}{Model Predictive Control}
\large

If trajectory optimization is fast enough, it can be used as a feedback controller.

\vspace{0.3cm}

MPC loop:

\begin{enumerate}
    \item Measure current state $\bm{x}(t)$.
    \item Solve a finite-horizon trajectory optimization problem.
    \item Execute only the first action.
    \item Move one step forward and repeat.
\end{enumerate}

\vspace{0.4cm}

\[
\pi(\bm{x})
=
\text{first action from the optimized trajectory}
\]

\vspace{0.3cm}

\begin{block}{Key idea}
MPC turns trajectory optimization into feedback by repeatedly re-planning.
\end{block}

\end{frame}

% Slide 18
\begin{frame}[t]{iLQR: big picture}
\large

iLQR solves nonlinear trajectory optimization by repeatedly turning the problem into a local time-varying LQR problem.

\vspace{0.25cm}

\begin{enumerate}
    \item Initialize a nominal control sequence:
    \[
    \bar{\bm{u}}_{0:N-1}
    \]
    \item Roll out the true dynamics:
    \[
    \bar{\bm{x}}_{t+1}=f_t(\bar{\bm{x}}_t,\bar{\bm{u}}_t)
    \]
    \item Around $(\bar{\bm{x}}_t,\bar{\bm{u}}_t)$, linearize dynamics and approximate cost quadratically.
    \item Solve this local LQR problem to get a better trajectory.
    \item Repeat until the true nonlinear cost no longer improves much.
\end{enumerate}

\vspace{0.15cm}

\[
\text{rollout}
\rightarrow
\text{local LQR}
\rightarrow
\text{improved rollout}
\rightarrow
\text{repeat}
\]

\end{frame}

% Slide 19
\begin{frame}[t]{iLQR problem setup}
\large

Discrete nonlinear dynamics:

\[
\bm{x}_{t+1}=f_t(\bm{x}_t,\bm{u}_t),
\qquad
\bm{x}_0 \text{ given}
\]

Finite-horizon objective:

\[
\min_{\bm{u}_0,\ldots,\bm{u}_{N-1}}
J
=
\ell_f(\bm{x}_N)
+
\sum_{t=0}^{N-1}\ell_t(\bm{x}_t,\bm{u}_t)
\]

\vspace{0.2cm}

At each iteration, search near the nominal trajectory:

\[
\bm{x}_t=\bar{\bm{x}}_t+\delta\bm{x}_t,
\qquad
\bm{u}_t=\bar{\bm{u}}_t+\delta\bm{u}_t
\]

\begin{block}{Important}
The nominal trajectory is feasible because it is generated by rolling out the dynamics.
\end{block}

\end{frame}

% Slide 20
\begin{frame}[t]{Step 1a: local dynamics}
\large

Linearize the dynamics around each nominal point:

\[
f_t(\bar{\bm{x}}_t+\delta\bm{x}_t,\bar{\bm{u}}_t+\delta\bm{u}_t)
\approx
f_t(\bar{\bm{x}}_t,\bar{\bm{u}}_t)
+
A_t\delta\bm{x}_t
+
B_t\delta\bm{u}_t
\]

\[
A_t=
\left.\frac{\partial f_t}{\partial \bm{x}}\right|_{\bar{\bm{x}}_t,\bar{\bm{u}}_t},
\qquad
B_t=
\left.\frac{\partial f_t}{\partial \bm{u}}\right|_{\bar{\bm{x}}_t,\bar{\bm{u}}_t}
\]

Since $\bar{\bm{x}}_{t+1}=f_t(\bar{\bm{x}}_t,\bar{\bm{u}}_t)$:

\[
\delta\bm{x}_{t+1}
\approx
A_t\delta\bm{x}_t+B_t\delta\bm{u}_t
\]

\end{frame}

% Slide 21
\begin{frame}[t]{Step 1b: local quadratic cost}
\small

Approximate the running cost to second order around
$(\bar{\bm{x}}_t,\bar{\bm{u}}_t)$:

\[
\begin{aligned}
\ell_t(\bm{x}_t,\bm{u}_t)
\approx
&\ \ell_t
+
\ell_x^T\delta\bm{x}_t
+
\ell_u^T\delta\bm{u}_t \\
&+
\frac{1}{2}\delta\bm{x}_t^T\ell_{xx}\delta\bm{x}_t
+
\delta\bm{u}_t^T\ell_{ux}\delta\bm{x}_t
+
\frac{1}{2}\delta\bm{u}_t^T\ell_{uu}\delta\bm{u}_t .
\end{aligned}
\]

All derivatives are evaluated at the nominal point:

\[
\ell_x=
\left.\frac{\partial \ell_t}{\partial \bm{x}}\right|_{\bar{\bm{x}}_t,\bar{\bm{u}}_t},
\qquad
\ell_u=
\left.\frac{\partial \ell_t}{\partial \bm{u}}\right|_{\bar{\bm{x}}_t,\bar{\bm{u}}_t}
\]

Terminal cost:

\[
\ell_f(\bm{x}_N)
\approx
\ell_f(\bar{\bm{x}}_N)
+
\ell_{f,x}^T\delta\bm{x}_N
+
\frac{1}{2}\delta\bm{x}_N^T\ell_{f,xx}\delta\bm{x}_N
\]

\end{frame}

% Slide 22
\begin{frame}[t]{Step 2: backward pass}
\small

Assume the next value function is locally quadratic:

\[
V_{t+1}
\approx
V_0'
+
{V_x'}^T\delta\bm{x}_{t+1}
+
\frac{1}{2}\delta\bm{x}_{t+1}^T V_{xx}'\delta\bm{x}_{t+1}
\]

Define the local one-step Q-function:

\[
Q_t(\delta\bm{x}_t,\delta\bm{u}_t)
=
\ell_t(\delta\bm{x}_t,\delta\bm{u}_t)
+
V_{t+1}(A_t\delta\bm{x}_t+B_t\delta\bm{u}_t)
\]

For fixed $\delta\bm{x}_t$, solve the local control problem:

\[
\delta\bm{u}_t^*
=
\operatorname*{arg\,min}_{\delta\bm{u}_t}
Q_t(\delta\bm{x}_t,\delta\bm{u}_t)
\]

Because $Q_t$ is quadratic, the solution is an affine correction:

\[
\delta\bm{u}_t^*
=
k_t+K_t\delta\bm{x}_t,
\qquad
k_t=-Q_{uu}^{-1}Q_u,
\qquad
K_t=-Q_{uu}^{-1}Q_{ux}
\]

\vspace{0.1cm}

\[
\text{feedforward step } k_t
\quad + \quad
\text{feedback correction } K_t\delta\bm{x}_t
\]

\end{frame}

% Slide 23
\begin{frame}[t]{Step 3: forward pass}
\large

The backward pass solves only a local approximation.
To update the real trajectory, roll out the nonlinear dynamics again.

\vspace{0.2cm}

For a step size $\alpha\in(0,1]$:

\[
\bm{x}^{\mathrm{new}}_0=\bm{x}_0
\]

\[
\bm{u}^{\mathrm{new}}_t
=
\bar{\bm{u}}_t
+
\alpha k_t
+
K_t(\bm{x}^{\mathrm{new}}_t-\bar{\bm{x}}_t)
\]

\[
\bm{x}^{\mathrm{new}}_{t+1}
=
f_t(\bm{x}^{\mathrm{new}}_t,\bm{u}^{\mathrm{new}}_t)
\]

\vspace{0.15cm}

Accept the first rollout that reduces the true nonlinear cost:

\[
J(\bm{x}^{\mathrm{new}},\bm{u}^{\mathrm{new}})
<
J(\bar{\bm{x}},\bar{\bm{u}})
\]

\end{frame}

% Slide 24
\begin{frame}[t]{Full iLQR loop}
\small

\begin{enumerate}
    \item Choose an initial control guess $\bar{\bm{u}}_{0:N-1}$.
    \item Roll out the dynamics to get $\bar{\bm{x}}_{0:N}$ and cost $J$.
    \item Backward pass:
    \begin{itemize}
        \item compute $A_t,B_t$ and cost derivatives,
        \item compute $Q_x,Q_u,Q_{xx},Q_{ux},Q_{uu}$,
        \item solve for $k_t$ and $K_t$.
    \end{itemize}
    \item Forward pass:
    \begin{itemize}
        \item try $\alpha=1,0.5,0.25,\ldots$,
        \item roll out true dynamics,
        \item accept the first lower-cost trajectory.
    \end{itemize}
    \item Replace the nominal trajectory and repeat.
\end{enumerate}

\begin{block}{Practical note}
If $Q_{uu}$ is not positive definite enough, use regularization
$Q_{uu}^{\mathrm{reg}}=Q_{uu}+\lambda I$.
\end{block}

\end{frame}


\end{document}
